% Part: lambda-calculus
% Chapter: introduction
% Section: lambda-definability

\documentclass[../../../include/open-logic-section]{subfiles}

\begin{document}

\olfileid{lam}{int}{rep}
\olsection{الدوال الحسابية \usetoken{S}{lambda definable}}

كيف يمكن أن يكون حساب لامبدا نموذجًا للحساب؟ لا يتضح في البداية حتى
كيف نفهم هذا القول. فللكلام على قابلية الحساب على الأعداد الطبيعية،
نحتاج إلى تمثيل مناسب لهذه الأعداد. والتمثيل الآتي ناجح على نحو مدهش.

\begin{defn}
لكل عدد طبيعي~$n$، عرّف \emph{عدد تشيرش}~$\num{n}$ بأنه حد لامبدا
$\lambd[x][\lambd[y][(x(x(x(\dots x(y)))))]]$، حيث يوجد في المجموع
$n$ من رموز~$x$.
\end{defn}

الحدود~$\num{n}$ «مكررات»: فعند إدخال~$f$، يعيد~$\num{n}$ الدالة التي
تربط~$y$ بـ$f^n(y)$. ولاحظ أن كل عدد سوي. ويمكننا الآن أن نقول ما معنى
أن «يحسب» حد لامبدا دالة على الأعداد الطبيعية.

\begin{defn}
لتكن~$f(x_0, \dots, x_{k-1})$ دالة جزئية $k$-موضعية من~$\Nat$ إلى
$\Nat$. نقول إن حد~$\lambd$~$F$ \emph{!!{lambda define}s}~$f$ إذا وفقط
إذا كان، لكل متتالية من الأعداد الطبيعية~$n_0$، \dots،~$n_{k-1}$،
\[
F\, \num{n_0}\, \num{n_1} \dots \num{n_{k-1}} \red \num{f(n_0, n_1, \dots,
  n_{k-1})}
\]
متى كانت~$f(n_0, \dots, n_{k-1})$ معرّفة، ولم يكن للحد
$F\, \num{n_0}\, \num{n_1} \dots \num{n_{k-1}}$ صورة سوية في غير ذلك.
\end{defn}

\begin{thm}
\ollabel{thm:lambda-def}
تكون الدالة~$f$ دالة جزئية قابلة للحساب إذا وفقط إذا كانت
!!{lambda defined} بحد لامبدا.
\end{thm}

\begin{explain}
هذه المبرهنة لافتة بعض الشيء. فحساب لامبدا، بوصفه نموذجًا للحساب، حساب
بسيط جدًا؛ وعمليتاه الوحيدتان هما تجريد لامبدا والتطبيق!{} ومع ذلك،
يمكن بهذه الموارد القليلة تنفيذ أي إجراء حسابي.
\end{explain}

\end{document}
